%
% 52-meromorph.tex -- Das Argumentprinzip für beliebige meromorphe Funktionen
%
% (c) 2026 Prof Dr Andreas Müller
%
\subsection{Das Argumentprinzip für beliebige meromorphe Funktionen}
Der Satz~\ref{buch:analytisch:argument:satz:ordnung}
liefert Information für eine einzelne isolierte Singularität.
Wir möchten jetzt zeigen, dass sich die Schlussfolgerung auch
verallgemeinern lässt auf einen Weg, der mehr als eine Null-
oder Polstelle der Funktion $f$ umschliesst.

\begin{satz}[Argumentprintzip für meromorphe Funktionen]
\label{buch:analytisch:argument:satz:argument}
\index{Argumentprinzip!meromorphe Funktionen}%
Sei jetzt $f$ ein meromorphe Funktion und $\gamma$ ein einfach
geschlossener Weg in $\mathbb{C}$, der keine Nullstellen oder Polstellen
von $f$ trifft.
Dann gilt
\[
\frac{1}{2\pi i}
\oint_\gamma
\frac{f'(z)}{f(z)}
\,dz
=
\sum_{\text{$a$ Singularität von $f$}}
\omega(a;f),
\]
wobei die Summe über alle Null- oder Polstellen $a$ von $f$ 
zu erstrecken ist.
\end{satz}

\begin{proof}
Nach Satz~\ref{buch:integration:homologie:satz:pfadpunkte} ist das
Wegintegral über $\gamma$ gleich gross wie die Summe der Integrale
über kleine, im Gegenuhrzeigersinn durchlaufene Kreise um die einzelnen
Singularitäten.
Die Summanden haben nach
Satz~\ref{buch:analytisch:argument:satz:ordnung}
den Wert $\omega(a;f)$.
Damit ist der Satz bewiesen.
\end{proof}
